\documentclass[UTF8,a4paper,12pt,fontset=windows]{ctexart}
\usepackage{xurl}
\urlstyle{same}

\usepackage[left=2.45cm,right=2.45cm,top=2.45cm,bottom=2.45cm,headsep=0.7cm,footskip=1.0cm]{geometry}
\usepackage{fontspec}
\usepackage{graphicx}
\usepackage{booktabs}
\usepackage{longtable}
\usepackage{tabularx}
\usepackage{array}
\usepackage{float}
\usepackage{hyperref}
\usepackage[table]{xcolor}
\usepackage{enumitem}
\usepackage{fancyhdr}
\usepackage{titlesec}
\usepackage{lastpage}
\usepackage{caption}
\usepackage{ragged2e}
\usepackage{setspace}
\usepackage{tocloft}
\usepackage{listings}

\setmainfont{Times New Roman}
\setsansfont{Times New Roman}
\setmonofont{Times New Roman}
\setCJKmainfont[BoldFont=SimHei]{SimSun}
\setCJKsansfont{Microsoft YaHei}
\setCJKmonofont{FangSong}

\definecolor{RuleGray}{RGB}{96,96,96}
\definecolor{TableHead}{RGB}{238,242,247}
\definecolor{TableAlt}{RGB}{250,251,253}

\hypersetup{
    hidelinks,
    pdftitle={YH_rv_cpu 初赛技术说明书},
    pdfauthor={CICC1003618}
}

\setstretch{1.50}
\setlength{\parindent}{2em}
\setlength{\parskip}{0.15em}
\setlength{\headheight}{15pt}
\emergencystretch=2em
\renewcommand{\arraystretch}{1.18}
\setlist[itemize]{leftmargin=2.2em,itemsep=0.25em,topsep=0.35em}
\setlist[enumerate]{leftmargin=2.4em,itemsep=0.25em,topsep=0.35em}

\titleformat{\section}{\centering\heiti\zihao{3}}{\thesection}{0.8em}{}
\titleformat{\subsection}{\raggedright\heiti\zihao{4}}{\thesubsection}{0.7em}{}
\titleformat{\subsubsection}{\raggedright\heiti\zihao{-4}}{\thesubsubsection}{0.7em}{}
\titlespacing*{\section}{0pt}{1.2ex plus 0.2ex minus 0.2ex}{0.8ex}
\titlespacing*{\subsection}{0pt}{0.9ex plus 0.2ex minus 0.1ex}{0.55ex}
\titlespacing*{\subsubsection}{0pt}{0.7ex plus 0.1ex minus 0.1ex}{0.4ex}

\captionsetup{font=small,labelfont=bf,labelsep=space,justification=centering}
\renewcommand{\contentsname}{目\quad 录}
\renewcommand{\cfttoctitlefont}{\hfill\heiti\zihao{3}}
\renewcommand{\cftaftertoctitle}{\hfill}
\renewcommand{\cftsecfont}{\normalsize}
\renewcommand{\cftsecpagefont}{\normalsize}
\setlength{\cftbeforesecskip}{2pt}

\newcolumntype{Y}{>{\RaggedRight\arraybackslash}X}
\newcolumntype{L}[1]{>{\RaggedRight\arraybackslash}p{#1}}

\lstdefinestyle{reportcode}{
    basicstyle=\ttfamily\small,
    frame=single,
    rulecolor=\color{RuleGray},
    backgroundcolor=\color{black!2},
    columns=fullflexible,
    breaklines=true,
    keepspaces=true,
    showstringspaces=false
}
\lstset{style=reportcode}

\newcommand{\Code}[1]{\texttt{#1}}
\newcommand{\Metric}[1]{{\rmfamily\addfontfeatures{Numbers={Lining,Proportional}}#1}}
\newcommand{\TeamID}{CICC1003618}
\newcommand{\TeamName}{芯银河战队}
\newcommand{\ProjectName}{YH\_rv\_cpu}
\newcommand{\ReportTitle}{YH\_rv\_cpu 初赛技术说明书}
\newcommand{\ReportDate}{2026-04-30}
\newcommand{\BoardName}{Xilinx PYNQ-Z2}
\newcommand{\BoardPart}{xc7z020clg400-1}
\newcommand{\BoardClock}{50.0\,MHz}
\newcommand{\CoreMarkScore}{4.137461}
\newcommand{\CoreMarkIter}{413.746136}
\newcommand{\CoreMarkTicks}{2416941}
\newcommand{\DhrystoneDMIPSMHz}{2.908287}
\newcommand{\BoardLUT}{4961}
\newcommand{\BoardFF}{2367}
\newcommand{\BoardBRAM}{6}
\newcommand{\BoardDSP}{15}
\newcommand{\BoardWNS}{+0.338\,ns}
\newcommand{\BoardWHS}{+0.059\,ns}
\newcommand{\CoverLogoPath}{./logo.png}
\newcommand{\BitstreamName}{YH\_rv\_cpu\_pynq\_z2\_sysclk\_8p000ns\_cpu50.bit}

\pagestyle{fancy}
\fancyhf{}
\lhead{\ProjectName}
\rhead{初赛技术说明书}
\cfoot{\thepage}


\newcommand{\CoverField}[2]{
\noindent
\makebox[0.16\textwidth][l]{\heiti\zihao{4} #1}
\underline{\makebox[0.70\textwidth][l]{\rule{0pt}{2.6ex}\zihao{-4} #2}}
\par\vspace{1.0em}
}

\newcommand{\FigureAsset}[2]{
\begin{figure}[H]
\centering
\IfFileExists{figures/pngs/#1}{
    \includegraphics[width=0.94\textwidth]{figures/pngs/#1}
}{
    \fcolorbox{RuleGray}{black!2}{
    \begin{minipage}[c][5.0cm][c]{0.90\textwidth}
    \centering
    \small
    \textbf{待插入图片}\\[0.4em]
    \Code{figures/pngs/#1}
    \end{minipage}
    }
}
\caption{#2}
\end{figure}
}

\begin{document}

\begin{titlepage}
    \centering
    \vspace*{0.4cm}
    \IfFileExists{\CoverLogoPath}{
        \includegraphics[width=6.2cm]{\CoverLogoPath}\par
    }{
        \fbox{\begin{minipage}[c][3.4cm][c]{6.2cm}\centering\heiti 大赛 LOGO\end{minipage}}\par
    }
    \vspace{0.6cm}
    {\heiti\zihao{2} 第十届\par}
    \vspace{0.7cm}
    {\heiti\zihao{2} 全国大学生集成电路创新创业大赛\par}
    \vfill
    \begin{minipage}{0.86\textwidth}
        \CoverField{报告类型：}{技术说明书}
        \CoverField{参赛杯赛：}{“七星微”企业命题}
        \CoverField{作品名称：}{基于 RISC-V 的五级流水 CPU 设计与 FPGA 验证}
        \CoverField{队伍编号：}{\TeamID}
        \CoverField{团队名称：}{\TeamName}
    \end{minipage}
\end{titlepage}

\thispagestyle{empty}


\clearpage
\pagestyle{empty}
\tableofcontents
\clearpage

\setcounter{page}{1}
\pagestyle{fancy}



\clearpage
\pagenumbering{arabic}
\setcounter{page}{1}
\pagestyle{fancy}
\section*{重要内容预览}
\addcontentsline{toc}{section}{重要内容预览}

\begin{table}[H]
\centering
\caption{作品核心亮点与实测结果}
\rowcolors{2}{TableAlt}{white}
\begin{tabularx}{\textwidth}{L{3.1cm}Y}
\toprule
\rowcolor{TableHead}
项目 & 结果 \\
\midrule
CPU 架构 & RV32I 五级顺序流水，包含 IF、ID、EX、MEM、WB，配套最小 SoC 与 FPGA 顶层。 \\
指令与优化 & 最终提交路径支持 RV32I、Zmmul、Zba/Zbb/Zbs 与 \Code{JAL} 早跳转；Zbc、XThead indexed memidx、IDBR 作为参数化探索路径，保留回归测试记录。 \\
CoreMark 结果 & \Metric{\CoreMarkScore{} CoreMark/MHz}，\Metric{\CoreMarkIter{} iter/s @ 100MHz}，\Metric{\CoreMarkTicks{} ticks}。 \\
FPGA 结果 & \BoardName{}，CPU \BoardClock{}，\Metric{\BoardLUT{} LUT / \BoardFF{} FF / \BoardBRAM{} BRAM / \BoardDSP{} DSP}，满足 \Metric{LUT < 5000}。 \\
时序与上板 & 实现后 \Metric{WNS=\BoardWNS{}}、\Metric{WHS=\BoardWHS{}}，硬件下载日志检出 \Code{PROGRAM\_OK: xc7z020\_1}。 \\
\bottomrule
\end{tabularx}
\end{table}

本文面向初赛评审，围绕架构设计、RTL 建模、验证结果与 FPGA 适配四条主线展开，重点呈现本作品在高性能、低资源占用与工程可复现性方面的实现亮点。
\clearpage

\section{作品概述}

\subsection{设计目标}

\ProjectName{} 是面向 FPGA 原型验证的轻量级 RISC-V CPU。设计目标是在 \Metric{LUT < 5000} 的资源约束内，完成可综合、可仿真、可上板的五级流水处理器原型，并兼顾性能提升、控制复杂度和低功耗导向实现。为此，设计重点保留对基准程序最敏感的热点路径与热点指令扩展，避免引入面积和时序代价较高但收益有限的复杂结构。

工程实现包括：
\begin{itemize}
    \item ISA：\Code{RV32I + Zmmul + Zba/Zbb/Zbs}，并保留 Zbc、XThead、IDBR 参数化探索路径；
    \item 控制流优化：\Code{JAL early redirect} 与低面积静态分支处理；
    \item FPGA：\BoardName{}，器件 \BoardPart{}，CPU \BoardClock{}；
    \item 性能指标：\Metric{\CoreMarkScore{} CoreMark/MHz}。
\end{itemize}

\subsection{设计依据与方法}

RISC-V 开放 ISA 具有模块化扩展特征，便于处理器在性能、面积和功耗之间按应用场景做取舍\cite{riscv-unprivileged,case-riscv}。本设计遵循“基线指令正确、热点扩展增强、控制路径受限前移”的方法：以 RV32I 五级流水作为可验证基线，以 Zmmul 和 Bitmanip 子扩展覆盖 CoreMark、Dhrystone 等整型程序中的乘法、地址生成、位操作与条件选择热点\cite{riscv-bitmanip}，同时避免引入会显著推高 LUT 与时序压力的复杂乱序结构。

在微结构层面，五级顺序流水仍是 FPGA 软核 CPU 中兼顾教学可解释性、时序收敛和资源效率的常用结构。相关 RISC-V 软核研究表明，面向分支、前递、访存和乘法路径做局部优化，可在保持轻量实现的同时取得可观性能收益\cite{rvcorep,rvcorep-m}。本作品据此将优化集中在可量化验证的短路径上：分支与 \Code{JAL} 的前端重定向、Zmmul 乘法保留、Bitmanip 热点指令保留，以及可关闭的参数化扩展组织。与应用级大核常见的高面积结构相比，单发射顺序流水更适合本赛题 \Metric{LUT < 5000} 的硬约束和 PYNQ-Z2 原型平台\cite{ariane-energy}。

\subsection{核心设计亮点}

本设计围绕轻量级高性能 RISC-V 处理器的常见优化方向展开，重点落在前端控制流缩短、热点扩展保留和 PPA 受约束条件下的选择性增强。
\begin{itemize}
    \item \textbf{热点扩展选择性保留：}冻结路径保留 \Code{Zmmul} 与 \Code{Zba/Zbb/Zbs}，面向乘法、位操作和地址生成等高频模式提升吞吐；Zbc 与 XThead 子集作为参数化探索路径，用于展示进一步性能潜力；
    \item \textbf{前端控制流提前化：}在不显著拉长关键路径的前提下，把部分分支比较与 \Code{JAL} 跳转前移至 ID 阶段，降低控制冒险的恢复开销；
    \item \textbf{低资源低功耗导向：}采用五级顺序流水、同步存储体与参数化扩展开关，移除硬件除法器与高代价控制逻辑，降低面积、切换活动和时序压力；
    \item \textbf{工程闭环完整：}形成从 RTL、TestBench、基准程序、综合实现到上板下载日志的可追溯证据链，便于复现与评审核查。
\end{itemize}

\section{CPU 架构设计}

\subsection{总体结构}

系统采用“CPU Core + 最小 SoC + FPGA 顶层”的三层组织方式。CPU Core 负责指令取指、译码、执行、访存、写回、CSR 和异常处理；SoC 层连接同步指令存储、数据存储与 MMIO；FPGA 顶层完成 PYNQ-Z2 时钟、复位、按键、LED 和 UART 适配。

\FigureAsset{01-system-architecture.png}{系统总体架构图}

\subsection{五级流水线}

处理器核心为五级顺序流水结构，阶段划分如下：
\begin{table}[H]
\centering
\caption{流水线阶段职责}
\rowcolors{2}{TableAlt}{white}
\begin{tabularx}{\textwidth}{L{2.4cm}Y}
\toprule
\rowcolor{TableHead}
阶段 & 职责 \\
\midrule
IF & 维护 PC，发起指令存储访问，处理 redirect、预测跳转和复位入口。 \\
ID & 译码指令，读取寄存器，生成立即数、控制信号、分支早判定与冒险请求。 \\
EX & 执行 ALU、乘法、位操作、分支/JAL/JALR 目标计算和异常条件判断。 \\
MEM & 处理 load/store、字节使能、非对齐检查和同步数据存储访问。 \\
WB & 选择 ALU/MEM/PC+4/CSR 结果写回寄存器堆。 \\
\bottomrule
\end{tabularx}
\end{table}

\FigureAsset{02-pipeline-control.png}{五级流水与关键控制通路图}

\subsection{指令集与性能优化}

面向 CoreMark、Dhrystone 与通用整型程序热点，处理器重点保留以下低成本高收益优化：
\begin{itemize}
    \item \textbf{Zmmul：}保留硬件乘法并移除硬件除法路径，使乘法操作吞吐量得益于专用乘法器，同时控制 LUT 与关键路径增长；
    \item \textbf{Zba/Zbb/Zbs：}增强地址计算、逻辑融合与位提取能力，提升热点位操作与地址生成效率，并在 PYNQ-Z2 上保持 \Metric{LUT < 5000}；
    \item \textbf{Zbc/XThead/IDBR 探索路径：}相关 RTL 与定向回归测试作为性能探索材料予以保留，但在最终冻结的 bitstream 中，为降低 LUT 资源占用与时序压力而未启用；
    \item \textbf{JAL early redirect：}对 \Code{JAL rd!=x0} 在 ID 阶段提前给出目标 PC，并在 EX 阶段完成一致性校验，降低前端重定向开销。
\end{itemize}

\subsection{控制冒险处理}

控制冒险统一通过 redirect 机制处理。EX 阶段保留分支、JAL、JALR、trap、mret 的完整重定向能力；ID 阶段仅对跳转目标可提前确定的路径实施重定向。基准测试统计显示，引入 \Code{JAL} 早跳转后，EX 阶段 \Code{JAL} redirect 计数降为 \Metric{0}，CoreMark 总完成周期由 \Metric{1974371} 降至 \Metric{1971888}，说明前端重定向代价得到有效压缩。

\section{Verilog 建模规范}

\subsection{模块边界}

工程采用 Verilog/SystemVerilog 建模，核心 RTL 目录按功能拆分：
\begin{itemize}
    \item \Code{YH\_rv\_cpu.v}：流水线顶层、控制流、冒险控制、级间寄存器；
    \item \Code{YH\_rv\_cpu\_id\_stage.v} 与 \Code{YH\_rv\_cpu\_decoder.v}：译码与控制信号生成；
    \item \Code{YH\_rv\_cpu\_ex\_stage.v} 与 \Code{YH\_rv\_cpu\_alu.v}：执行、乘法和位操作；
    \item \Code{YH\_rv\_cpu\_mem\_stage.v}：访存宽度、写掩码和非对齐检查；
    \item \Code{YH\_rv\_cpu\_soc.v}：ROM、RAM、MMIO 和仿真/FPGA 共用 SoC。
\end{itemize}

\subsection{编码约定}

\begin{enumerate}
    \item 参数化开关集中在顶层和 SoC 边界传递，例如 \Code{ENABLE\_ZMMUL\_EXTENSION}、\Code{ENABLE\_BITMANIP\_EXTENSION}、\Code{ENABLE\_ID\_BRANCH\_EX\_FORWARD}。
    \item 时序逻辑使用 \Code{always @(posedge clk or negedge rst\_n)}，组合逻辑使用 \Code{always @*} 或连续赋值，避免隐式 latch。
    \item 复位值在级间寄存器中显式给出，flush、stall 与 trap 路径优先级在 CPU 顶层集中处理。
    \item 新增指令优先通过 decoder、ALU opcode 和定向 TestBench 三处闭环，避免只改执行单元而遗漏非法指令路径。
    \item FPGA 可选功能必须能通过参数关闭，以便在 LUT 预算与性能之间快速切换。
\end{enumerate}

\subsection{注释与可追溯性}

源代码注释重点说明参数用途、控制流优先级、异常处理路径、扩展指令译码以及 TestBench 判定条件。提交源码包保留构建脚本、仿真脚本、测试程序和摘要日志，便于评委从文档指标追溯到原始工程证据。

\section{验证计划与测试报告}

\subsection{验证策略}

验证分为四层：指令与模块定向测试、SoC smoke 测试、性能基准测试、FPGA 综合实现与下载测试。文中所有性能与实现指标均来自脚本化构建和日志解析，不依赖人工填报数据。

\FigureAsset{03-validation-closure.png}{验证闭环与证据流图}

\begin{table}[H]
\centering
\caption{回归测试矩阵}
\rowcolors{2}{TableAlt}{white}
\begin{tabularx}{\textwidth}{L{3.2cm}Y L{2.4cm}}
\toprule
\rowcolor{TableHead}
测试项 & 覆盖内容 & 结果 \\
\midrule
SoC smoke & 启动、ROM/RAM、UART 输出、基本指令流。 & PASS \\
Zmmul 诊断 & \Code{mul} 正确性与 \Code{divu} 非支持异常路径。 & PASS \\
Bitmanip 提交路径 & Zba、Zbb、Zbs 快速子集和非支持 Zbc 拒绝路径。 & PASS \\
探索扩展路径 & Zbc、XThead condmov/ext/addsl、indexed memidx 与 IDBR 参数化回归。最终冻结 bitstream 未启用。 & PASS/未启用 \\
Branch predict diag & 静态 backward \Code{BNE} 预测与控制流回退。 & PASS \\
CoreMark & \Metric{10} iterations，2K performance profile，CRC 匹配。 & \Metric{\CoreMarkScore{} CoreMark/MHz} \\
Dhrystone & Dhrystone 2.2，运行 \Metric{10} 次，采用优化构建并复验结果。 & \Metric{\DhrystoneDMIPSMHz{} DMIPS/MHz} \\
PYNQ-Z2 实现 & 综合、布局布线、bitstream、硬件下载。 & PASS \\
\bottomrule
\end{tabularx}
\end{table}

\subsection{关键性能结果}

\begin{table}[H]
\centering
\caption{性能测试结果}
\rowcolors{2}{TableAlt}{white}
\begin{tabularx}{\textwidth}{L{3.4cm}Y}
\toprule
\rowcolor{TableHead}
指标 & 结果 \\
\midrule
CoreMark/MHz & \Metric{\CoreMarkScore{}} \\
CoreMark ticks & \Metric{\CoreMarkTicks{}} \\
Iter/s @ 100MHz & \Metric{\CoreMarkIter{}} \\
Dhrystone DMIPS/MHz & \Metric{\DhrystoneDMIPSMHz{}} \\
指令覆盖说明 & 覆盖 RV32I 基础指令、Zmmul、Zba/Zbb/Zbs 及非法/非支持指令路径；Zbc、XThead、IDBR 探索路径保留定向回归记录。 \\
分支预测说明 & 定向诊断 PASS；Profile 显示 \Code{JAL} EX redirect 降为 \Metric{0}，相关错误事件未检出。 \\
\bottomrule
\end{tabularx}
\end{table}

\subsection{优化前后差异}

以未采用本作品控制流优化组合的 RV32IM 参考路径 \Metric{4.060276 CoreMark/MHz} 为对照，最终冻结路径达到 \Metric{\CoreMarkScore{} CoreMark/MHz}，并在 Dhrystone 优化构建复验中达到 \Metric{\DhrystoneDMIPSMHz{} DMIPS/MHz}。性能提升主要来自 Zmmul 与 Bitmanip 热点扩展、低面积前端控制流优化和可关闭参数化组织；与此同时，设计保持 \Metric{\BoardLUT{} LUT} 的资源占用与正裕量时序结果，满足 \Metric{LUT < 5000} 的板级约束。高性能探索路径曾达到更高 CoreMark，但因 \Metric{6147 LUT} 与 \Metric{WNS=-0.801 ns} 未进入冻结提交口径。

\section{FPGA 适配指南}

\subsection{板级配置}

\begin{table}[H]
\centering
\caption{PYNQ-Z2 适配配置}
\rowcolors{2}{TableAlt}{white}
\begin{tabularx}{\textwidth}{L{3.2cm}Y}
\toprule
\rowcolor{TableHead}
项目 & 配置 \\
\midrule
开发板 & \BoardName{} \\
器件 & \BoardPart{} \\
输入时钟 & PYNQ-Z2 板载 \Metric{125 MHz} PL 时钟 \\
CPU 时钟 & MMCM 生成 \BoardClock{}，满足不低于 \Metric{50 MHz} 要求 \\
约束文件 & \Code{fpga\_artifacts\_pynq\_z2/constraints/pynq\_z2\_template.xdc} \\
bitstream & PYNQ-Z2 提交 bitstream，完整文件名见证据索引。 \\
下载证据 & Vivado Hardware Manager 日志，含 \Code{PROGRAM\_OK: xc7z020\_1}。 \\
\bottomrule
\end{tabularx}
\end{table}

\subsection{资源与时序}

\begin{table}[H]
\centering
\caption{实现后资源与时序结果}
\rowcolors{2}{TableAlt}{white}
\begin{tabularx}{\textwidth}{L{3.2cm}Y}
\toprule
\rowcolor{TableHead}
指标 & 结果 \\
\midrule
Slice LUTs & \Metric{\BoardLUT{} / 53200 = 9.33\%} \\
Slice Registers & \Metric{\BoardFF{} / 106400 = 2.22\%} \\
Block RAM Tile & \Metric{\BoardBRAM{} / 140 = 4.29\%} \\
DSP & \Metric{\BoardDSP{} / 220 = 6.82\%} \\
Setup WNS & \Metric{\BoardWNS{}}，0 个失败端点 \\
Hold WHS & \Metric{\BoardWHS{}}，0 个失败端点 \\
硬件下载 & Vivado Hardware Manager 检出 \Code{xc7z020\_1} 并完成 \Code{PROGRAM\_OK} \\
\bottomrule
\end{tabularx}
\end{table}

\FigureAsset{04-fpga-bringup-flow.png}{FPGA 原型适配路径图}

\subsection{上板连接说明}

JTAG 下载和器件识别通过 PYNQ-Z2 Micro-USB 完成。软核 UART 约束到 Pmod B，若需采集串口日志，应使用外接 \Metric{3.3 V} USB-UART 模块连接：
\begin{itemize}
    \item \Code{uart\_rxd\_out}：FPGA TX，Pmod B1 / Y14；
    \item \Code{uart\_txd\_in}：FPGA RX，Pmod B0 / W14；
    \item 电平：\Metric{3.3 V}。
\end{itemize}
板级验证证据包括 bitstream 下载日志、器件识别、实现报告、仿真回归结果、基准测试摘要和 PL 侧 UART 文本输出。PYNQ-Z2 的 Micro-USB 串口连接至 PS 端，不直接连接 PL 软核 UART；本作品将 PL 侧 UART 约束到 Pmod B，并使用外接 \Metric{3.3 V} CP2102 USB-UART 模块在 \Metric{115200 8N1} 参数下采集到连续 \Code{YH\_rv\_cpu CoreMark/MHz=4.137461 DMIPS/MHz=2.908287 tick=XX pc=XXXXXXXX} 输出。其中 \Code{tick} 为板级实时计数，\Code{pc} 为 CPU 调试 PC 采样，可作为 bitstream 下载后 FPGA 原型系统持续运行的上板证据。

\section{AI 工具使用声明}

本项目在非核心创作环节使用 AI 工具辅助研发与材料整理，工具包括 OpenAI 与 DeepSeek。AI 未替代核心 CPU 架构决策、RTL 最终确认和实验结论判定。

\begin{table}[H]
\centering
\caption{AI 工具使用情况}
\rowcolors{2}{TableAlt}{white}
\begin{tabularx}{\textwidth}{L{2.5cm}Y L{2.6cm}}
\toprule
\rowcolor{TableHead}
工具 & 使用场景 & 使用比例估计 \\
\midrule
OpenAI & 工程脚本梳理、CoreMark 热点分析、优化方案筛选、回归命令组织、日志快速核对。 & 文档约 20--30\%；代码修改建议低于 5\% \\
DeepSeek & 技术文稿润色、表述压缩、章节逻辑检查。 & 文档润色约 10--20\% \\
\bottomrule
\end{tabularx}
\end{table}

\section{结论}

\ProjectName{} 形成了完整的初赛设计证据链：CPU 架构与 RTL 建模边界明确，功能验证与性能测试可追溯，PYNQ-Z2 下载成功，资源占用与时序结果满足板级约束。该设计在保持轻量级实现和低资源占用的同时，通过热点扩展保留、前端控制流优化和参数化工程组织取得了 \Metric{\CoreMarkScore{} CoreMark/MHz} 与 \Metric{\DhrystoneDMIPSMHz{} DMIPS/MHz} 的实测表现，体现了在 FPGA 资源受限条件下对性能、面积与工程可实现性的综合平衡。

\clearpage
\addcontentsline{toc}{section}{参考文献}
\small
\begin{thebibliography}{9}
\bibitem{riscv-unprivileged}
RISC-V International. The RISC-V Instruction Set Manual, Volume I: Unprivileged Architecture. \url{https://riscv.github.io/riscv-isa-manual/snapshot/unprivileged/}

\bibitem{riscv-bitmanip}
RISC-V International. Bit-Manipulation Extension. \url{https://docs.riscv.org/reference/isa/unpriv/b-st-ext.html}

\bibitem{case-riscv}
K. Asanovic and D. A. Patterson. Instruction Sets Should Be Free: The Case For RISC-V. EECS Technical Report, 2014. \url{https://people.eecs.berkeley.edu/~krste/papers/EECS-2014-146.pdf}

\bibitem{rvcorep}
T. Miyazaki et al. RVCoreP: An Optimized RISC-V Soft Processor of Five-Stage Pipelining. IEICE Transactions on Information and Systems, 2020. \url{https://arxiv.org/abs/2002.03568}

\bibitem{rvcorep-m}
T. Miyazaki et al. RVCoreP-32IM: An Effective Architecture to Implement Mul/Div Instructions for Five Stage RISC-V Soft Processors. 2020. \url{https://arxiv.org/abs/2010.16171}

\bibitem{ariane-energy}
F. Zaruba et al. The Cost of Application-Class Processing: Energy and Performance Analysis of a Linux-ready 1.7GHz 64-bit RISC-V Core in 22nm FD-SOI Technology. 2019. \url{https://arxiv.org/abs/1904.05442}
\end{thebibliography}

\clearpage
\section*{证据索引}
\addcontentsline{toc}{section}{证据索引}

\begin{tabularx}{\textwidth}{L{4.3cm}Y}
\toprule
\rowcolor{TableHead}
证据项 & 路径 \\
\midrule
CoreMark 摘要 & \Code{性能与验证报告/性能日志/} 中的提交路径 CoreMark 摘要。 \\
Dhrystone 摘要 & \Code{性能与验证报告/性能日志/} 中的 Dhrystone 摘要。 \\
实现资源报告 & \Code{FPGA原型系统/fpga\_artifacts\_pynq\_z2/reports/impl\_utilization.rpt} \\
实现时序报告 & \Code{FPGA原型系统/fpga\_artifacts\_pynq\_z2/reports/impl\_timing\_summary.rpt} \\
硬件下载日志 & \Code{FPGA原型系统/fpga\_artifacts\_pynq\_z2/logs/} 中的 PYNQ-Z2 下载日志。 \\
bitstream & \Code{FPGA原型系统/fpga\_artifacts\_pynq\_z2/} 中的 PYNQ-Z2 bitstream。 \\
\bottomrule
\end{tabularx}

\end{document}
